\chapter{What would make the bank better off?}
\label{reading:banking-decision}

\section{More business is not necessarily more value}

Consider a deposit offer to a defined group of households. The bank can send
the offer or leave its current terms unchanged. For household $i$, let
$Y_i(1)$ be the contribution over the next quarter if the offer is sent and
$Y_i(0)$ the contribution under the alternative. Contribution includes the
agreed funding value, interest expense and servicing costs over the same
quarter. The average effect for the eligible population is
\[
 \tau=\mathbb E\{Y_i(1)-Y_i(0)\}.
\]
The expectation averages over households, rather than accounts or branches.
That choice prevents a household with several accounts from receiving several
votes merely because of the way the records are stored. The definition would
change if management wanted the average effect per pound of existing deposits.
The formal potential-outcome development follows in Chapter~\ref{r3:graphs}.

A simple financial calculation illustrates why the outcome needs care. Suppose
a rate increase attracts an additional balance $\Delta B$ for one year. Let
$m$ be the contribution per pound of that extra balance before the cost of the
rate increase, $\Delta r$ the increase in the annual rate, $B_0$ the existing
balance that also receives it, and $c$ the direct cost of the offer. Under the
illustrative assumptions that balances remain constant during the year and
that the increase applies to both existing and additional balances, net value is
\begin{align}
 V&=m\Delta B-\Delta r(B_0+\Delta B)-c\\
  &=(m-\Delta r)\Delta B-\Delta r B_0-c.
 \label{reading:deposit-value}
\end{align}
The first term in the second line is the contribution from additional funding
after repricing. The second is the cost of paying more on balances the bank
already had. The last is the direct cost. A large balance response can still
leave $V$ negative. The example omits changes in credit risk, capital needs,
competitor responses and balances within the year; a business case that depends
on those quantities must include them in the outcome definition.

For illustration, take $B_0=\pounds10{,}000$, $m=0.03$, $\Delta r=0.0025$
and $c=\pounds5$. The bank needs more than
\[
 \Delta B^*=\frac{0.0025\times10{,}000+5}{0.03-0.0025}
 =\pounds1{,}090.91
\]
of additional balance to break even. An increase of \pounds1,000 gives
$V=-\pounds2.50$; \pounds2,000 gives $V=\pounds25$. These are stipulated
calculations, not observations or forecasts. They show what a causal estimate
must measure before it can inform the financial decision.

\begin{figure}[htbp]\centering
\begin{tikzpicture}[x=2.25cm,y=0.055cm,>=Stealth]
 \draw[->] (0,-35)--(0,66) node[above,align=center]{\small Net annual value\ (\pounds)};
 \draw[->] (0,-30)--(3.35,-30) node[right,align=left]{\small Additional balance\ (\pounds000)};
 \draw[dashed,gray] (0,0)--(3.1,0);
 \draw[very thick,teal] (0,-30)--(3,52.5);
 \foreach \x in {1,2,3}{\draw (\x,-32)--(\x,-28);\node[below] at (\x,-32){\small\x};}
 \foreach \y in {0,25,50}{\draw (-.035,\y)--(.035,\y);\node[left] at (-.04,\y){\small\y};}
 \fill[ink] (1.09091,0) circle (2pt);
 \node[above left,align=right,font=\small] at (1.09091,0){Break-even balance\\\pounds1,091};
\end{tikzpicture}
\caption{A deposit campaign must earn back the cost of repricing existing
balances. Illustrative calculation from equation~\ref{reading:deposit-value};
the line is a specified financial relation, not a fitted demand curve.}
\label{reading:break-even}
\end{figure}

\section{Selection can imitate an effect}

Comparing customers who accepted an offer with those who did not generally
mixes the offer's effect with the reasons for acceptance. The distinction can
be seen without a complicated model. Suppose, as a constructed example, that
half the eligible customers have high underlying demand for deposits and half
have low demand. Their balances without an offer would be \pounds200 and
\pounds100 in the units of the example. Assume the offer adds \pounds20 for
every customer. If 80\% of recipients are high-demand customers but only 20\%
of nonrecipients are, observed average balances are
\begin{align*}
 \mathbb E[Y\mid A=1]&=0.8(220)+0.2(120)=200,\\
 \mathbb E[Y\mid A=0]&=0.2(200)+0.8(100)=120.
\end{align*}
The observed difference is \pounds80, although the specified effect is
\pounds20. The remaining \pounds60 is a difference in the composition of the
two groups. The calculation assumes equally sized recipient and nonrecipient
groups, consistent with the population shares. Random assignment in this
example would make the groups comparable in expectation; an observational
analysis needs an additional argument about demand and treatment selection.

\begin{figure}[htbp]\centering
\begin{tikzpicture}[>=Stealth, every node/.style={align=center,font=\small}]
 \node[draw,rounded corners,fill=light,text width=3.0cm] (d) at (0,1.8)
   {Underlying demand\ and financial circumstances};
 \node[draw,rounded corners,fill=light,text width=2.8cm] (a) at (-2.7,0)
   {Receiving\ the deposit offer};
 \node[draw,rounded corners,fill=light,text width=2.8cm] (y) at (2.7,0)
   {Subsequent\ deposit balance};
 \draw[->,thick] (d)--(a);\draw[->,thick] (d)--(y);
 \draw[->,very thick,teal] (a)--(y) node[midway,below]{Effect of interest};
\end{tikzpicture}
\caption{A proposed explanation of selection into an offer. The upper path
creates an association even without the causal effect shown in teal. The
arrows are assumptions about the constructed example.}
\label{reading:selection}
\end{figure}

This is why identification precedes precision. A more accurate estimate of
\pounds80 would not answer the stated question. Nor does predicting balances
well reveal which part of the difference the bank can change. Randomized
invitations, eligibility rules, sufficiently rich pre-offer information and
credible structural restrictions offer different ways forward. They identify
different quantities under different assumptions, developed in
Chapters~\ref{r3:graphs} and~\ref{r3:investment} and the intervening analysis.

\section{Timing determines what a decision can use}

Information available after a customer accepts an offer cannot ordinarily
justify a decision made before acceptance. The distinction is easy to miss
when every field appears in one table. An observed call duration, for example,
belongs to a completed call. It cannot be used as though it were known when
deciding whether to make that call. Section~\ref{pd15:datasets} examines the
roles and limitations of public financial data in more detail.

\begin{figure}[htbp]\centering
\begin{tikzpicture}[>=Stealth,every node/.style={font=\small,align=center}]
 \draw[->,thick] (0,0)--(11.8,0);
 \foreach \x in {0.8,3.8,6.8,10.2}{\draw (\x,-.12)--(\x,.12);}
 \node[above,text width=2.2cm] at (.8,.25){Recorded income,\ balances and history};
 \node[below,text width=2.2cm] at (3.8,-.25){Offer\ decision};
 \node[above,text width=2.2cm] at (6.8,.25){Acceptance and\ subsequent behaviour};
 \node[below,text width=2.5cm] at (10.2,-.25){Funding contribution\ at the stated horizon};
 \draw[dashed,teal] (3.8,-1.2)--(3.8,1.5);
\end{tikzpicture}
\caption{The information date separates predictors available to the decision
from later responses. This is an illustrative sequence, not a claim about a
particular bank's records.}
\label{reading:time}
\end{figure}

Repeated servicing decisions make the problem more demanding. A contact today
may change both tomorrow's repayment and tomorrow's information. Comparing
contacted and uncontacted borrowers without accounting for that history can
confuse a response to financial distress with a cause of it. A policy must
therefore specify what the bank observes when it acts. The later longitudinal
analysis compares such rules, including cases in which relevant borrower
information is hidden.

\section{Uncertainty has an economic consequence}

A confidence interval describes sampling uncertainty under stated assumptions.
It does not establish those assumptions. A set of effects compatible with
several causal models answers a different question: how much ambiguity remains
even if the observed relationships were known. Both can affect the decision.
Their difference is developed formally in the identification and inference
chapters; it should also remain visible in a management report.

\begin{figure}[htbp]\centering
\begin{tikzpicture}[x=.065cm,y=.9cm,>=Stealth,every node/.style={font=\small}]
 \draw[->] (-35,0)--(65,0) node[right]{Net value (\pounds)};
 \draw[dashed,gray] (0,0)--(0,3.5);
 \foreach \x in {-20,0,20,40,60}{\draw (\x,-.08)--(\x,.08);\node[below] at (\x,-.1){\x};}
 \foreach \y/\lo/\hi/\mid in {1/10/40/25,2/-10/40/15,3/-25/-5/-15}{
   \draw[very thick,teal] (\lo,\y)--(\hi,\y);
   \draw[thick] (\lo,\y-.12)--(\lo,\y+.12);
   \draw[thick] (\hi,\y-.12)--(\hi,\y+.12);
   \fill[ink] (\mid,\y) circle (2pt);
 }
 \node[left] at (-36,1){A};\node[left] at (-36,2){B};\node[left] at (-36,3){C};
\end{tikzpicture}
\caption{Three stipulated ranges for the same net-value decision. A lies above
break-even, B crosses it, and C lies below it. These are schematic uncertainty
ranges, not estimated confidence intervals. Their role is to explain why
precision matters through the choice it can change.}
\label{reading:uncertainty}
\end{figure}

If the plausible values all favour the same action, greater precision may be
worth little. When the preferred action changes across plausible values, better
information may be valuable. The relevant information could concern customer
identity, assignment to an offer, a missing hardship measure or the financial
cost of an intervention. It need not come from a larger sample of the same
records. The investment analysis returns to this distinction when comparing
the cost of another experiment with the decision it might improve.

\section{The case for a shared analytical service}

The proposal follows from this economic sequence. A common service could help
analysts retain the population, action, information date and financial outcome
as a study moves from formulation to estimation and interpretation. Existing
methods should be reused where their assumptions and capabilities fit. An
analysis that finds insufficient information should explain the missing
comparison and, where possible, the measurement or experiment that would
resolve it.

The value of language assistance must be assessed within that sequence. A
clearer formulation or a useful criticism of a causal diagram can save analyst
time. A fluent but unsupported assumption can make the study less reliable.
The comparisons in Part VI therefore concern particular research tasks and
their errors, rather than general impressions of intelligence. Public evidence
can test some of these tasks before confidential bank data are needed.

The scientific and investment questions meet at the same point: whether a
better-supported decision is worth its cost. The chapters that follow supply
the causal arguments, statistical methods, software comparisons and evidence
needed to make that judgment.
